\title{Controlling Text-to-Image Diffusion by Orthogonal Finetuning}

\begin{abstract}
Large-scale text-to-image diffusion models demonstrate remarkable abilities in creating photorealistic images from textual descriptions. Effectively steering or controlling these powerful models to perform various downstream tasks remains a significant challenge. To address this issue, we propose a principled finetuning approach called Orthogonal Finetuning~(OFT) for adapting text-to-image diffusion models to specific downstream applications. Unlike current methods, OFT can theoretically maintain the hyperspherical energy, which describes the pairwise relationships between neurons on the unit hypersphere. We observe that preserving this property is vital for maintaining the semantic generation capabilities of text-to-image diffusion models. To enhance finetuning stability, we introduce Constrained Orthogonal Finetuning (COFT), which adds an extra radius constraint to the hypersphere. In particular, we focus on two key finetuning tasks for text-to-image generation: subject-driven generation, where the objective is to produce subject-specific images based on a few images of a subject and a text prompt, and controllable generation, which aims to allow the model to incorporate additional control signals. Our empirical results demonstrate that the OFT framework surpasses existing approaches in both generation quality and convergence speed.
\end{abstract}

\section{Introduction}

Recent text-to-image diffusion models~\cite{saharia2022photorealistic,ramesh2022hierarchical,rombach2022high} have achieved remarkable success in text-guided high-fidelity image synthesis. Despite these advances, guidance based solely on text can remain ambiguous and insufficient for delivering detailed and precise control over generated images. To overcome this limitation, this work focuses on two types of text-to-image generation tasks:

\begin{itemize}[leftmargin=*,nosep]

    \item \textbf{Subject-driven generation}~\cite{ruiz2023dreambooth}: Given only a few images depicting a subject, the goal is to generate images of the same subject placed in novel contexts, guided by a text prompt.
    \item \textbf{Controllable generation}~\cite{zhang2023adding,mou2023t2i}: Given an extra control signal (\emph{e.g.}, canny edges, segmentation maps), the task is to create images that comply with this control input in addition to the text prompt.
\end{itemize}

Both tasks essentially reduce to the challenge of how to finetune text-to-image diffusion models effectively without compromising the generative performance acquired during pretraining. We outline the key requirements for an effective finetuning approach as: (1) \emph{training efficiency}: minimizing the number of trainable parameters and training epochs, and (2) \emph{generalizability preservation}: maintaining high-fidelity and diverse generative capabilities. Typically, finetuning is performed either by slightly adjusting neuron weights with a small learning rate (\emph{e.g.}, \cite{ruiz2023dreambooth}) or by introducing a minor component with re-parameterized neuron weights (\emph{e.g.}, \cite{hulora2022,zhang2023adding}). Although straightforward, neither of these finetuning strategies can ensure the preservation of pretrained generative performance. Moreover, there is a lack of a principled framework for designing an effective finetuning approach and selecting appropriate hyperparameters, such as the number of training epochs. A significant challenge lies in the absence of a metric to quantify the extent to which pretrained generative capabilities are preserved. Current finetuning techniques implicitly assume that a smaller Euclidean distance between the finetuned and pretrained models corresponds to better retention of pretrained abilities. For this reason, finetuning methods generally employ very small learning rates. While this assumption may hold in certain cases, we contend that Euclidean distance alone does not sufficiently capture the semantic preservation level. Hence, adopting a more structural metric to characterize the differences between the finetuned and pretrained models can substantially improve both the preservation of pretrained performance and the stability of finetuning.

Motivated by empirical findings that hyperspherical similarity effectively encodes semantic information~\cite{liu2017deep,liu2018decoupled,chen2020angular}, we utilize hyperspherical energy~\cite{liu2018learning} to describe the pairwise relational structure among neurons. Hyperspherical energy is defined as the sum of hyperspherical similarities (such as cosine similarity) between every pair of neurons within the same layer, reflecting the degree of neuron uniformity on the unit hypersphere~\cite{liu2021learning}. We posit that an effective finetuned model should exhibit minimal change in hyperspherical energy relative to the pretrained model. A straightforward approach involves adding a regularization term to maintain constant hyperspherical energy during finetuning; however, this does not guarantee effective minimization of the hyperspherical energy difference. To address this, we exploit an invariance property of hyperspherical energy—namely, the pairwise hyperspherical similarity is theoretically preserved if an identical orthogonal transformation is applied to all neurons. Building on this invariance, we introduce Orthogonal Finetuning (OFT), which adapts large text-to-image diffusion models for downstream tasks without altering their hyperspherical energy. The core concept is to learn a shared orthogonal transformation at the layer level that preserves the pairwise angles between neurons. OFT can also be interpreted as modifying the canonical coordinate system for neurons within the same layer. Further considering that smaller Euclidean distances between the finetuned and pretrained models correlate with better retention of pretraining performance, we propose a variant called Constrained Orthogonal Finetuning (COFT), which restricts the finetuned model to lie within a hypersphere of fixed radius centered on the pretrained neurons.

The rationale behind why orthogonal transformation is effective for finetuning neurons partially stems from the 2D Fourier transform, which allows an image to be broken down into magnitude and phase spectra. The phase spectrum, representing the angular relationship between the input and basis, retains the majority of the semantic content. For instance, an image’s phase spectrum combined with a random magnitude spectrum can still reconstruct the original image without losing its semantic meaning. This observation implies that altering the directions of neurons is fundamental to semantically modifying the generated image—an objective central to both subject-driven and controllable generation. Nonetheless, modifying neuron directions with excessive freedom will inevitably compromise the generative performance established during pretraining. To limit this freedom, we propose maintaining the angle between every pair of neurons, based largely on the hypothesis that these inter-neuron angles are vital for encoding the neural network's knowledge. Guided by this insight, it is natural to learn a layer-shared orthogonal transformation for neurons within each layer so that the hyperspherical energy remains constant.

We also take inspiration from orthogonal over-parameterized training~\cite{liu2021orthogonal}, which trains classification neural networks from scratch by applying orthogonal transformations to a randomly initialized network. This is motivated by the fact that a randomly initialized network exhibits provably low hyperspherical energy on average, and the objective in \cite{liu2021orthogonal} is to maintain this low hyperspherical energy throughout training (as lower energy correlates with better generalization in classification~\cite{liu2018learning,lin2020regularizing}). The work in \cite{liu2021orthogonal} demonstrates that orthogonal transformations provide sufficient flexibility to train neural networks that generalize well for classification tasks. By contrast, our focus lies in finetuning text-to-image diffusion models to enhance controllability and improve downstream generative performance. We highlight two key distinctions between OFT and \cite{liu2021orthogonal}. First, whereas \cite{liu2021orthogonal} aims to minimize hyperspherical energy, OFT is designed to preserve the pretrained model’s hyperspherical energy, thereby safeguarding the intrinsic pretrained structure during finetuning. Minimizing hyperspherical energy during diffusion model finetuning risks damaging the original semantic configurations. Second, OFT strives to minimize deviation from the pretrained model, which results in a constrained variant, unlike \cite{liu2021orthogonal} which does not impose such constraints. Achieving a suitable balance between flexibility and stability is crucial in finetuning, and we contend that our OFT framework successfully fulfills this objective. Our contributions are summarized as follows:

\begin{itemize}[leftmargin=*,nosep]

    \item We introduce a new finetuning technique called Orthogonal Finetuning, designed to enhance controllability in text-to-image diffusion models. To boost stability further, we propose a constrained version that restricts the angular deviation from the pretrained model.
    \item Unlike existing finetuning approaches, OFT finetunes the model while mathematically ensuring the preservation of hyperspherical energy, which we empirically determine to be a crucial indicator of the pretrained model’s generative semantic integrity.
    \item We apply OFT to two scenarios: subject-driven generation and controllable generation. Through an extensive empirical evaluation, we show that OFT significantly outperforms previous methods in generation quality, convergence speed, and finetuning robustness. Additionally, OFT demonstrates enhanced sample efficiency, achieving convergence with substantially fewer training images and epochs.
    \item For controllable generation, we propose a novel control consistency metric to assess controllability. The key idea involves estimating the control signal from the generated image and comparing it to the original control input. This metric further confirms OFT’s strong controllability.
\end{itemize}

\section{Related Work}

\textbf{Text-to-image diffusion models}. Considerable advances~\cite{saharia2022photorealistic,ramesh2022hierarchical,rombach2022high,gu2022vector,nichol2022glide} have been achieved in text-to-image generation, largely driven by progress in diffusion-based generative models~\cite{ho2020denoising,song2021score,song2021denoising,dhariwal2021diffusion} and vision-language representation learning~\cite{su2020vlbert,lu2019vilbert,radford2021learning,wang2021simvlm,li2022blip,li2023blip,alayrac2022flamingo,schuhmann2022laion}. GLIDE~\cite{nichol2022glide} and Imagen~\cite{saharia2022photorealistic} operate diffusion models directly in pixel space. GLIDE jointly trains the text encoder and a diffusion prior on paired text-image data, whereas Imagen employs a frozen pretrained text encoder. Stable Diffusion~\cite{rombach2022high} and DALL-E2~\cite{ramesh2022hierarchical} perform diffusion modeling in latent space. Stable Diffusion learns a latent space visual codebook via VQ-GAN~\cite{esser2021taming}, while DALL-E2 uses CLIP~\cite{radford2021learning} to form a shared latent embedding space for images and text. Beyond diffusion models, generative adversarial networks~\cite{reed2016generative,zhang2017stackgan,xu2018attngan,li2019controllable} and autoregressive models~\cite{ramesh2021zero,ding2021cogview,wu2022nuwa,yuscaling} have also been explored for text-to-image generation. OFT is inherently model-agnostic and applicable to any text-to-image diffusion model.

\textbf{Subject-driven generation}. To avoid altering the subject, \cite{avrahami2022blended,nichol2022glide} introduce user-provided masks as additional conditions. Inversion techniques~\cite{choi2021ilvr,dhariwal2021diffusion,ramesh2022hierarchical,gal2022image} allow contextual modifications without changing the subject. \cite{hertz2022prompt} enables both local and global edits without requiring input masks. However, these approaches struggle to maintain identity-specific details of the subject. Pivotal Tuning~\cite{roich2022pivotal} finetunes a generator around an initial inverted latent code with extra regularization to preserve identity. Similarly, \cite{nitzan2022mystyle} learns a personalized generative prior for faces from a collection of an individual’s images. \cite{casanova2021instance} produces diverse variations of an instance but may lose instance-specific features. DreamBooth~\cite{ruiz2023dreambooth} finetunes a text-to-image diffusion model with a customized token and a small set of subject images, using a reconstruction loss plus a class-specific prior preservation loss. OFT builds upon the DreamBooth framework, but replaces naive small-rate finetuning with orthogonal transformation-based finetuning.

\textbf{Controllable generation}. Image-to-image translation can be interpreted as a type of controllable generation, with prior approaches predominantly utilizing conditional generative adversarial networks~\cite{isola2017image,zhu2017toward,wang2018high,park2019semantic,choi2018stargan}. Diffusion models have also been employed for image-to-image translation tasks~\cite{saharia2022palette,voynov2022sketch,wang2022pretraining}. Recently, ControlNet~\cite{zhang2023adding} introduced a method to steer a pretrained diffusion model by finetuning it to incorporate extra control signals, achieving remarkable performance in controllable generation. Another contemporary and related study, T2I-Adapter~\cite{mou2023t2i}, similarly finetunes a pretrained diffusion model to enhance control over the generated images. Building on the same task framework as \cite{zhang2023adding,mou2023t2i}, we utilize OFT to finetune pretrained diffusion models, consistently delivering superior controllability with reduced training data requirements and fewer finetuning parameters. Crucially, OFT does not incur any additional computational cost during inference.

\textbf{Model finetuning}. The practice of finetuning large pretrained models for downstream tasks has grown increasingly widespread~\cite{devlin2018bert,brown2020language,he2022masked}. Adaptation techniques (e.g., \cite{hulora2022,houlsby2019parameter,pfeiffer2020adapterfusion}) represent a category of finetuning that has been extensively examined in natural language processing. Among these, LoRA~\cite{hulora2022} is the method most closely related to OFT, as it assumes a low-rank structure for the additive weight updates during finetuning. By contrast, OFT applies a layer-shared orthogonal transformation to modify neuron weights multiplicatively, which has been theoretically shown to preserve the pairwise angles between neurons within the same layer, leading to improved stability.

\section{Orthogonal Finetuning}

\subsection{Why Does Orthogonal Transformation Make Sense?}

We begin by exploring the rationale for employing orthogonal transformation in finetuning text-to-image diffusion models. This question is broken down into two sub-questions: (1) why it is beneficial to finetune the angles (i.e., directions) of neurons, and (2) why orthogonal transformation is chosen for adjusting these angles.

Regarding the first question, we draw motivation from the empirical findings in \cite{liu2018decoupled,chen2020angular} which indicate that angular feature differences effectively represent the semantic gap. SphereNet~\cite{liu2017deep} demonstrates that normalizing all neurons onto a unit hypersphere during neural network training preserves comparable capacity and even enhances generalizability, suggesting that neuron directions alone can encapsulate the critical information from the data. To further illustrate the significance of neuron angles, we conduct a simple experiment shown in Figure~\ref{fig:toy_ae}, where a standard convolutional autoencoder is trained on flower images. During training, the standard inner product is employed to generate the feature map ($z$ represents the output element of the convolution kernel $\bm{w}$, and $\bm{x}$ denotes the input within the sliding window). At test time, we compare three methods for producing the feature map: (a) the inner product used during training, (b) the magnitude information, and (c) the angular information. The results depicted in Figure~\ref{fig:toy_ae} reveal that angular information from neurons can nearly perfectly reconstruct the input images, whereas the magnitude information carries no meaningful content. We highlight that cosine activation (c) is not applied during training, which relies solely on the inner product. These findings suggest that neuron angles (directions) predominantly store the semantic information of input images. Consequently, to alter the semantic content of images, fine-tuning neuron directions is likely to be more effective.

Addressing the second question, the most straightforward approach to fine-tune neuron directions is by simultaneously rotating or reflecting all neurons within the same layer, naturally introducing orthogonal transformations. While it could be more flexible to employ other angular transformations that rotate neurons by different angles individually, orthogonal transformation strikes a balance between flexibility and structure. Additionally, \cite{liu2021orthogonal} confirms that orthogonal transformations possess sufficient capacity for neural network learning. To substantiate our claim, we conduct an experiment demonstrating the strong regularization effect imposed by the orthogonality constraint. We execute a controllable generation experiment following the ControlNet~\cite{zhang2023adding} framework, with results displayed in Figure~\ref{fig:ortho}. It is evident that our standard OFT maintains stable performance and achieves precise control after training completion (epoch 20). In contrast, OFT without the orthogonality constraint fails to produce realistic images and exhibits no control ability. This experiment validates the critical role of the orthogonality constraint in OFT.

\subsection{General Framework}

The traditional finetuning approach generally employs gradient descent with a small learning rate to update a model (or specific layers within a model). Using a small learning rate implicitly promotes only minor deviations from the pretrained model, and conventional finetuning essentially seeks to train the model while implicitly minimizing $\thickmuskip=2mu \medmuskip=2mu \| \bm{M} - \bm{M}^0\|$, where $\bm{M}$ represents the finetuned model weights and $\bm{M}^0$ denotes the pretrained model weights. This implicit regularization is intuitively reasonable; however, it can still permit excessive flexibility when finetuning a large model. To mitigate this, LoRA incorporates an additional low-rank constraint on the weight update, specifically, $\thickmuskip=2mu \medmuskip=2mu \text{rank}(\bm{M} - \bm{M}^0)=r'$, where $r'$ is chosen to be a small integer. In contrast to LoRA, OFT imposes a constraint on the pairwise neuron similarity: $\thickmuskip=2mu \medmuskip=2mu \| \text{HE}(\bm{M})-\text{HE}(\bm{M}^0) \|=0$, with $\text{HE}(\cdot)$ indicating the hyperspherical energy of a weight matrix. To illustrate, consider a fully connected layer $\thickmuskip=2mu \medmuskip=2mu \bm{W}=\{\bm{w}_1,\cdots,\bm{w}_n\}\in\mathbb{R}^{d\times n}$ where each neuron $\bm{w}_i \in \mathbb{R}^d$ (and $\bm{W}^0$ corresponds to the pretrained weights). The output vector $\thickmuskip=2mu \medmuskip=2mu \bm{z} \in \mathbb{R}^n$ of this layer is computed by $\thickmuskip=2mu \medmuskip=2mu \bm{z} = \bm{W}^\top \bm{x}$ where $\thickmuskip=2mu \medmuskip=2mu \bm{x} \in \mathbb{R}^d$ is the input vector. OFT can be viewed as minimizing the difference in hyperspherical energy between the finetuned and pretrained models:

\begin{equation}\label{eq:oft_principle}
\footnotesize
\min_{\bm{W}} \left\| \text{HE}(\bm{W})-\text{HE}(\bm{W}^0)\right\|~~\Leftrightarrow~~\min_{\bm{W}} \bigg{\|} \sum_{i\neq j} \|\hat{\bm{w}}_i-\hat{\bm{w}}_j\|^{-1}-\sum_{i\neq j} \|\hat{\bm{w}}^0_i-\hat{\bm{w}}^0_j\|^{-1}\bigg{\|}
\end{equation}

Here, $\thickmuskip=2mu \medmuskip=2mu \hat{\bm{w}}_i:=\bm{w}_i/\|\bm{w}_i\|$ denotes the normalized $i$-th neuron, and the hyperspherical energy of the fully connected layer $\bm{W}$ is defined as $\thickmuskip=2mu \medmuskip=2mu \text{HE}(\bm{W}):= \sum_{i\neq j} \|\hat{\bm{w}}_i - \hat{\bm{w}}_j\|^{-1}$. It is straightforward to see that the minimum achievable value for Eq.~\eqref{eq:oft_principle} is zero. This minimum is attained whenever $\bm{W}$ differs from $\bm{W}^0$ solely by a rotation or reflection, i.e., $\thickmuskip=2mu \medmuskip=2mu \bm{W} = \bm{R} \bm{W}^0$, where $\thickmuskip=2mu \medmuskip=2mu \bm{R} \in \mathbb{R}^{d \times d}$ is an orthogonal matrix (with determinant $1$ or $-1$ indicating rotation or reflection, respectively). This embodies the fundamental concept of OFT: finetuning the network by learning layer-shared orthogonal matrices that transform neurons

\begin{equation}\label{eq:oft_general}
    \footnotesize
    \bm{z}=\bm{W}^\top\bm{x}=(\bm{R}\cdot\bm{W}^0)^\top\bm{x},~~~\text{s.t.}~\bm{R}^\top\bm{R}=\bm{R}\bm{R}^\top=\bm{I}
\end{equation}

Here, $\bm{W}$ represents the OFT-finetuned weight matrix and $\bm{I}$ is the identity matrix. OFT is depicted in Figure~\ref{fig:block}. Analogous to the zero initialization approach used in LoRA, it is necessary to ensure that OFT fine-tunes the pretrained model starting exactly from $\bm{W}^0$. To accomplish this, we initialize the orthogonal matrix $\bm{R}$ as the identity matrix so that the fine-tuned model begins with the pretrained weights. To maintain the orthogonality of $\bm{R}$, differential orthogonalization techniques described in \cite{liu2021orthogonal,lezcano2019cheap} can be applied. We will further explain how to efficiently preserve orthogonality from a computational perspective.

\subsection{Efficient Orthogonal Parameterization}\label{sect:eff_ortho}

Although conventional orthogonalization methods like the Gram-Schmidt process are differentiable, they tend to be computationally expensive in practice~\cite{liu2021orthogonal}. To improve efficiency, we employ the Cayley parameterization to generate the orthogonal matrix. Concretely, we define the orthogonal matrix as $\thickmuskip=2mu \medmuskip=2mu \bm{R}=(\bm{I}+\bm{Q})(\bm{I}-\bm{Q})^{-1}$, where $\bm{Q}$ is a skew-symmetric matrix fulfilling $\thickmuskip=2mu \medmuskip=2mu \bm{Q}=-\bm{Q}^\top$. This method’s efficiency comes with a minor limitation — Cayley parameterization can only produce orthogonal matrices with determinant $1$, which belong to the special orthogonal group. Fortunately, we observe that this constraint does not negatively impact performance in practice. Even using the Cayley transform, $\bm{R}$ can become parameter-inefficient when the dimension $d$ is large. To mitigate this, we propose representing $\bm{R}$ as a block-diagonal matrix composed of $r$ blocks, resulting in:

\begin{equation}
    \footnotesize
    \bm{R}=\text{diag}(\bm{R}_1,\bm{R}_2,\cdots,\bm{R}_r)=\begin{bmatrix}
    \bm{R}_1\in \text{O}\left(\frac{d}{r}\right) &  &\\
    &\ddots & \\
    & & \bm{R}_r\in \text{O}\left(\frac{d}{r}\right)
    \end{bmatrix}\in \text{O}(d)
\end{equation}

where $\text{O}(d)$ represents the orthogonal group of dimension $d$, and $\thickmuskip=2mu \medmuskip=2mu \bm{R}\in\mathbb{R}^{d\times d}$ and $\thickmuskip=2mu \medmuskip=2mu \bm{R}_i\in\mathbb{R}^{d/r\times d/r},\forall i$ are orthogonal matrices. When $\thickmuskip=2mu \medmuskip=2mu r=1$, the block-diagonal orthogonal matrix reduces to the conventional unconstrained orthogonal matrix. For an orthogonal matrix of size $\thickmuskip=2mu \medmuskip=2mu d \times d$, the number of parameters is $\thickmuskip=2mu \medmuskip=2mu d(d-1)/2$, leading to a computational complexity of $\mathcal{O}(d^2)$. For an orthogonal matrix with $r$ block-diagonal blocks, the parameter count is $\thickmuskip=2mu \medmuskip=2mu d(d/r - 1)/2$, resulting in complexity $\thickmuskip=2mu \medmuskip=2mu \mathcal{O}(d^2/r)$. Furthermore, one may optionally share the block matrices to further decrease the parameter count, i.e., $\thickmuskip=2mu \medmuskip=2mu \bm{R}_i = \bm{R}_j$ for all $i \neq j$. This brings down the parameter complexity to $\mathcal{O}(d^2 / r^2)$. Despite these parameter efficiency enhancements, it is important to note that the resulting matrix $\bm{R}$ remains orthogonal, ensuring that hyperspherical energy preservation is not compromised.

Next, we compare the parameter efficiency of OFT with LoRA. For LoRA with a low-rank parameter $r'$, the total number of trainable parameters is $\thickmuskip=2mu \medmuskip=2mu r'(d + n)$. Assuming both $r$ and $r'$ scale with neuron dimension $d$ (e.g., $\thickmuskip=2mu \medmuskip=2mu r = r' = \alpha d$, where $\thickmuskip=2mu \medmuskip=2mu 0 < \alpha \leq 1$ is some constant), the parameter complexity of LoRA becomes $\mathcal{O}(d^2 + dn)$, whereas OFT’s complexity is $\mathcal{O}(d)$. To illustrate this difference concretely, consider a weight matrix sized $\thickmuskip=2mu \medmuskip=2mu 128 \times 128$. With $r' = 8$, LoRA has $2,048$ trainable parameters, while OFT with $r = 8$ (without block sharing) has only $960$ trainable parameters.

\subsection{Constrained Orthogonal Finetuning}

We can further restrict the flexibility of the original OFT by confining the finetuned model to remain close to the pretrained model. Specifically, COFT implements the forward pass as:

\begin{equation}\label{eq:coft_general}
    \footnotesize
    \bm{z} = \bm{W}^\top \bm{x} = (\bm{R} \cdot \bm{W}^0)^\top \bm{x}, \quad \text{s.t.} \quad \bm{R}^\top \bm{R} = \bm{R} \bm{R}^\top = \bm{I}, \quad \norm{\bm{R} - \bm{I}} \leq \epsilon
\end{equation}

which enforces both an orthogonality constraint and a bounded deviation $\epsilon$ from the identity matrix. The orthogonality constraint can be enforced using the Cayley parameterization described in Section~\ref{sect:eff_ortho}. However, integrating the $\epsilon$-deviation constraint into the Cayley parameterization is nontrivial. To better understand the Cayley transform, we employ the Neumann series to approximate $\thickmuskip=2mu \medmuskip=2mu \bm{R} = (\bm{I} + \bm{Q})(\bm{I} - \bm{Q})^{-1}$ as $\thickmuskip=2mu \medmuskip=2mu \bm{R} \approx \bm{I

series converges under the operator norm). Consequently, we can incorporate the constraint $\thickmuskip=2mu \medmuskip=2mu\|\bm{R}-\bm{I}\|\leq\epsilon$ directly within the Cayley transform, leading to an equivalent constraint of $\thickmuskip=2mu \medmuskip=2mu \|\bm{Q}-\bm{0}\|\leq\epsilon'$, where $\bm{0}$ represents a zero matrix and $\epsilon'$ is a distinct error hyperparameter (different from $\epsilon$). This new constraint on matrix $\bm{Q}$ can be straightforwardly enforced using projected gradient descent. To initialize the orthogonal matrix $\bm{R}$ as an identity matrix, we set $\bm{Q}$ to be the zero matrix initially. COFT can be regarded as integrating two explicit constraints: minimizing the hyperspherical energy difference and limiting the deviation from the pretrained model. The latter constraint is typically implicit in existing finetuning approaches, but COFT explicitly enforces it. Despite the strong performance of OFT, we find that COFT offers even greater stability in finetuning due to this explicit deviation constraint. Figure~\ref{fig:eps_coft} illustrates how $\epsilon$ impacts COFT's performance. It demonstrates that $\epsilon$ regulates finetuning flexibility: larger $\epsilon$ values make the COFT-finetuned model resemble the OFT-finetuned one, while smaller $\epsilon$ values cause the COFT-finetuned model to behave more like the original pretrained text-to-image diffusion model.

\subsection{Re-scaled Orthogonal Finetuning}

We introduce a straightforward extension to the original OFT by additionally learning a magnitude scaling factor for each neuron. This extension is motivated by the fact that re-scaling neurons does not affect the hyperspherical energy, as magnitudes are normalized out. Specifically, the forward pass is defined as $\thickmuskip=2mu \medmuskip=2mu\bm{z}=(\bm{R}\bm{W}^0\bm{D})^\top\bm{x}$\footnote{\textbf{Errata}: In the NeurIPS camera-ready version, the forward pass of re-scaled OFT was incorrectly written as $\thickmuskip=2mu \medmuskip=2mu\bm{z}=(\bm{D}\bm{R}\bm{W}^0)^\top\bm{x}$. The original implementation is correct, so the results reported in Appendix~\ref{app:rescaled} remain valid.} where $\thickmuskip=2mu \medmuskip=2mu \bm{D}=\text{diag}(s_1,\cdots,s_n)\in\mathbb{R}^{n\times n}$ is a learnable diagonal matrix with strictly positive diagonal entries $s_1,\cdots,s_n$. Unlike the original OFT forward pass in Eq.~\eqref{eq:oft_general}, where only $\bm{R}$ is learnable, here both the diagonal matrix $\bm{D}$ and the orthogonal matrix $\bm{R}$ are trainable parameters. This re-scaled OFT increases OFT’s flexibility with a negligible number of extra parameters. We adhere to the original OFT in our experiments to demonstrate the efficacy of orthogonal transformation alone, although we observe that the re-scaled OFT generally performs better (see Appendix~\ref{app:rescaled}).

\section{Intriguing Insights and Discussions}

\textbf{OFT is architecture-agnostic}. In principle, OFT can be applied to any neural network architecture. For Transformers, LoRA is commonly applied to the attention weights~\cite{hulora2022}. To ensure a fair comparison with LoRA, we restrict OFT to finetuning attention weights in our experiments. Beyond fully connected layers, OFT is also highly suitable for finetuning convolutional layers, since the block-diagonal structure of $\bm{R}$ provides meaningful interpretations in convolutional contexts (unlike LoRA). When the number of blocks matches the number of input channels, each block transforms a distinct neuron channel independently, resembling the learning of depth-wise convolution kernels~\cite{chollet2017xception}. Conversely, when all blocks in $\bm{R}$ are shared, OFT applies an orthogonal transformation to neurons that is shared across channels. We present preliminary experiments on finetuning convolution layers with OFT in Appendix~\ref{app:convolution}.

\textbf{Relation to LoRA}. LoRA adds a low-rank matrix to prevent significant changes in the pretrained weight matrix’s information. In contrast, OFT enforces the transformation applied to the pretrained weight matrix to be orthogonal (full-rank), thereby preserving the integrity of the pretraining information. The forward pass of OFT can be rewritten as $\thickmuskip=2mu \medmuskip=2mu \bm{z}=(\bm{R}\bm{W}^0)^\top\bm{x}=(\bm{W}^0+(\bm{R}-\bm{I})\bm{W}^0)^\top\bm{x}$, where $\thickmuskip=2mu \medmuskip=2mu (\bm{R}-\bm{I})\bm{W}^0$ serves a similar role to LoRA’s low-rank weight update. Given that $\bm{W}^0$ is generally full-rank, OFT also performs a low-rank weight update when $\thickmuskip=2mu \medmuskip=2mu \bm{R}-\bm{I}$ is low-rank. Analogous to LoRA’s rank parameter $r'$, OFT employs a diagonal block parameter $r$ to limit the number of trainable parameters. Notably, LoRA and OFT embody two different approaches to parameter efficiency: LoRA leverages low-rank structure to reduce trainable parameters, while OFT utilizes sparsity structure (i.e., block-diagonal orthogonality) to achieve efficiency.

\textbf{Reasons for OFT’s faster convergence}. First, as shown in Figure~\ref{fig:toy_ae}, the most effective way to modify semantic information is by altering neuron directions, which is precisely what OFT targets. Additionally, OFT can be interpreted as fine-tuning neurons constrained on a smooth hypersphere manifold, resulting in a more favorable optimization landscape. This advantage is also supported by empirical evidence in \cite{liu2021orthogonal}.

\textbf{Why we do not minimize hyperspherical energy}. Unlike \cite{liu2021orthogonal}, our approach does not aim to minimize hyperspherical energy. In classification tasks, it is preferable to have neurons without redundancy. Minimizing hyperspherical energy leads to neurons being evenly distributed around the hypersphere, which is not a desirable goal in fine-tuning since it risks erasing the knowledge gained during pretraining.

\textbf{Trade-off between flexibility and regularity in finetuning}. We identify a fundamental trade-off between flexibility and regularization. Standard finetuning offers the greatest flexibility but tends to suffer from poor stability and is prone to model collapse. Surprisingly simple, OFT achieves a favorable balance between flexibility and regularity by maintaining the pairwise angles between neurons. The block-diagonal parameterization can also be interpreted as a stronger form of regularization applied to the orthogonal matrix.

\textbf{No additional inference overhead}. In contrast to ControlNet, our OFT framework does not impose any extra inference cost on the finetuned model. During inference, we simply multiply the learned orthogonal matrix $\bm{R}$ with the pretrained weight matrix $\bm{W}^0$ to produce an equivalent weight matrix $\thickmuskip=2mu \medmuskip=2mu \bm{W}=\bm{R}\bm{W}^0$. Therefore, the inference speed remains identical to that of the pretrained model.

\section{Experiments and Results}

\textbf{General settings}. In our experiments, we use Stable Diffusion v1.5~\cite{rombach2022high} as the pretrained text-to-image model. To maintain fairness, images are randomly selected from each method’s outputs. For subject-driven generation, we broadly follow the setup of DreamBooth~\cite{ruiz2023dreambooth}. For controllable generation, our approach aligns generally with ControlNet~\cite{zhang2023adding} and T2I-Adapter~\cite{mou2023t2i}. To ensure an equitable comparison with LoRA, we apply either OFT or COFT exclusively to the same layer targeted by LoRA. Additional results and detailed information are provided in Appendix~\ref{app:settings}.

\subsection{Subject-driven Generation}

\textbf{Settings}. We use DreamBooth~\cite{ruiz2023dreambooth} and LoRA~\cite{hulora2022} as baseline methods. All approaches employ the same loss function as in DreamBooth. For DreamBooth and LoRA, we typically adhere to the original papers’ recommended hyperparameter settings. Further results are available in Appendix~\ref{app:settings}, \ref{app:coft_oft}, \ref{app:more_qualitative}, \ref{app:failure}.

\textbf{Finetuning stability and convergence}. We begin by assessing the finetuning stability and convergence speed of DreamBooth, LoRA, OFT, and COFT. The outcomes are presented in Figure~\ref{fig:teaser} and Figure~\ref{fig:db_convergence}. It is evident that both COFT and OFT can finetune the diffusion model with considerable stability. After 400 iterations, both DreamBooth and OFT variants demonstrate effective control, whereas LoRA struggles to maintain the subject’s identity. At 2000 iterations, DreamBooth starts producing collapsed images, and LoRA fails to generate the yellow shirt, instead producing yellow fur. Conversely, OFT and COFT continue to provide stable and consistent control over the generated images. These findings confirm the rapid yet stable convergence of our OFT framework in subject-driven generation. We highlight that the robustness to the number of finetuning iterations is crucial, as it alleviates the challenge of tuning iteration counts for different subjects. For both OFT and COFT, one can set a relatively large number of iterations directly without meticulous tuning. Moreover, with an appropriate $\epsilon$, COFT makes setting both the learning rate and iteration count straightforward.

\textbf{Quantitative comparison}. Following \cite{ruiz2023dreambooth}, we perform a quantitative comparison evaluating subject fidelity (DINO~\cite{caron2021emerging}, CLIP-I~\cite{radford2021learning}), text prompt fidelity (CLIP-T~\cite{radford2021learning}), and sample diversity (LPIPS~\cite{zhang2018unreasonable}). CLIP-I measures the average pairwise cosine similarity of CLIP embeddings between generated and real images. DINO is analogous to CLIP-I but uses ViT S/16 DINO embeddings. CLIP-T computes the average cosine similarity between the text prompt and generated image embeddings. Additionally, we assess average LPIPS cosine similarity among generated images of the same subject with the identical text prompt. Table~\ref{table:dreambooth_final} indicates that both COFT and OFT significantly outperform DreamBooth and LoRA on the DINO and CLIP-I metrics, while showing slightly improved or comparable results in prompt fidelity and diversity. For each approach, we finetune the same pretrained model 30 times with different random seeds to reduce randomness. The results demonstrate that our OFT framework not only ensures superior convergence and stability but also consistently achieves better final performance.

\textbf{Qualitative comparison}. To better illustrate the advantages of OFT, we present several randomly selected examples of subject-driven generation in Figure~\ref{fig:dreambooth_final}. For fairness, all examples are generated using the same finetuned model for each method, ensuring that no individual text prompt is separately optimized for the final outputs. For each approach, we choose the model that attains the highest validation CLIP metrics. As shown in Figure~\ref{fig:dreambooth_final}, both OFT and COFT demonstrate excellent semantic preservation of the subject, whereas LoRA frequently fails to maintain the subject's identity (for instance, LoRA completely loses the subject's identity in the bowl example). Additionally, OFT and COFT exhibit much more precise control through text prompts, while DreamBooth, despite preserving subject identity, often does not generate images that accurately follow the text prompt (e.g., the first row in the bowl example). This qualitative comparison highlights that our OFT framework simultaneously achieves superior controllability and subject preservation. Furthermore, OFT's performance is robust to the number of iterations, maintaining effectiveness even with a large iteration count, unlike DreamBooth or LoRA. Additional qualitative examples are provided in Appendix~\ref{app:more_qualitative}. We also include a human evaluation in Appendix~\ref{app:human} that further confirms the advantages of OFT.

\subsection{Controllable Generation}

\textbf{Settings}. We select ControlNet~\cite{zhang2023adding}, T2I-Adapter~\cite{mou2023t2i}, and LoRA~\cite{hulora2022} as baseline methods. The main paper focuses on three challenging controllable generation tasks: Canny edge to image (C2I) on the COCO dataset~\cite{lin2014microsoft}, segmentation map to image (S2I) on the ADE20K dataset~\cite{zhou2017scene}, and landmark to face (L2F) on the CelebA-HQ dataset~\cite{karras2018progressive,xia2021tedigan}. All methods are employed to finetune Stable Diffusion (SD) v1.5 on these datasets over 20 epochs. Further results can be found in Appendices~\ref{app:more_qualitative}, \ref{app:more_control_task}, and \ref{app:failure}.

\textbf{Convergence}. We assess the convergence speed of ControlNet, T2I-Adapter, LoRA, and COFT on the L2F task through both quantitative and qualitative analyses. For the quantitative measure, we calculate the average $\ell_2$ distance between the control face landmarks and the predicted face landmarks. In Figure~\ref{fig:face_convergence}, we display the face landmark error obtained from models fine-tuned over different numbers of epochs. It is evident that both COFT and OFT converge considerably faster. LoRA requires 20 epochs to reach the performance level that our OFT framework attains by the 8th epoch. We highlight that OFT and COFT have a comparable number of trainable parameters to LoRA (which is significantly fewer than ControlNet) while demonstrating much higher convergence efficiency compared to existing methods. Additionally, the rapid convergence of OFT is further supported by the results shown in Figure~\ref{fig:teaser}. The example on the right in Figure~\ref{fig:teaser} illustrates that OFT is substantially more data-efficient than both ControlNet and LoRA, as it can achieve good convergence using only 5\% of the ADE20K dataset. For qualitative comparisons, we concentrate on OFT, COFT, and ControlNet, since ControlNet achieves landmark errors closest to ours. The results in Figure~\ref{fig:face_convergence_qual} demonstrate that OFT and COFT both converge steadily, with the generated face pose progressively aligning with the control landmarks. In contrast, ControlNet shows a sudden onset of controllability after the 8th epoch, reflecting the abrupt convergence pattern noted in \cite{zhang2023adding}. This stability in convergence makes the OFT framework more practical, as its training dynamics are notably smoother and more predictable. Consequently, tuning OFT's hyperparameters is easier to accomplish.

\textbf{Quantitative comparison}. We propose a control consistency metric to assess the effectiveness of controllable generation. The core concept is to extract the control signal from the generated image and then compare it with the original input control signal. For the C2I task, we calculate the IoU and F1 score. For the S2I task, we measure mean IoU, mean accuracy, and overall accuracy. For the L2F task, we determine the average $\ell_2$ distance between the control landmarks and the predicted landmarks. Additional details about the consistency metrics are provided in Appendix~\ref{app:settings}. For all the methods compared, we utilize their optimal hyperparameter configurations. The results presented in Table~\ref{table:control_gen} demonstrate that both OFT and COFT achieve significantly stronger and more precise control than the other approaches. We note that adapter-based methods (\emph{e.g.}, T2I-Adapter and ControlNet) tend to converge more slowly and produce inferior final outcomes. Compared to ControlNet, LoRA shows superior performance in the S2I task but underperforms in the C2I and L2F tasks. Overall, we observe that LoRA's performance ceiling is relatively limited, even after careful hyperparameter tuning. In contrast, the performance of our OFT framework appears not to have reached saturation, as we empirically observe continued improvement with an increasing number of trainable parameters. We highlight that our quantitative evaluation method for controllable generation is a novel contribution, as it provides an accurate assessment of the control capabilities of fine-tuned models (subject to the accuracy of the off-the-shelf segmentation/detection model).

\textbf{Qualitative comparison}. We further conduct a qualitative comparison between OFT and COFT against leading state-of-the-art approaches such as ControlNet, T2I-Adapter, and LoRA. The randomly generated images presented in Figure~\ref{fig:control_final} indicate that OFT and COFT not only produce high-quality, realistic images but also maintain precise control over the outputs. For the S2I task, it is evident that LoRA fails to generate images that adhere to the input segmentation map, whereas ControlNet, OFT, and COFT effectively regulate the generated images. Compared to ControlNet, both OFT and COFT generate images with higher fidelity, richer details, and more plausible geometric structures, all while using significantly fewer model parameters. In the C2I task, OFT and COFT can create semantically consistent images from approximate Canny edges, whereas T2I-Adapter and LoRA show considerably poorer performance. For the L2F task, our approach delivers the most precise pose control in the generated faces, even under challenging facial orientations. Across all three control tasks, our results demonstrate that OFT and COFT produce qualitatively superior images relative to the current state-of-the-art baselines, underscoring the effectiveness of the OFT framework for controllable image generation. To offer a more extensive qualitative assessment, additional example results for all three control tasks are included in Appendix~\ref{app:control_exp}. Furthermore, we showcase the strong performance of OFT on additional control tasks such as dense pose to human body, sketch to image, and depth to image in Appendix~\ref{app:more_control_task}.

\section{Concluding Remarks and Open Problems}

Inspired by the insight that angular relationships among neurons play a key role in defining visual semantics, we introduce a straightforward yet powerful finetuning technique—orthogonal finetuning (OFT)—for controlling text-to-image diffusion models. Our approach focuses on two main text-to-image tasks: subject-driven generation and controllable generation. In comparison to existing methods, OFT offers enhanced controllability and more stable finetuning while using fewer finetuning parameters. Crucially, OFT does not add any inference-time overhead, resulting in an efficiently deployable model.

OFT also brings up several intriguing open questions. Firstly, OFT ensures orthogonality through Cayley parametrization, which requires computing a matrix inverse. This requirement somewhat constrains the scalability of OFT. Although we mitigate this issue by employing a block diagonal parametrization, accelerating the matrix inverse computation in a differentiable manner remains a challenging problem. Secondly, OFT holds unique promise in compositionality because the orthogonal matrices generated by multiple OFT finetuning tasks can be multiplied together and still yield an orthogonal matrix. Investigating whether this set of orthogonal matrices can retain knowledge from all downstream tasks is a compelling avenue for future research. Lastly, the parameter efficiency of OFT heavily relies on the block diagonal structure, which unavoidably introduces some bias and restricts flexibility. Finding ways to enhance parameter efficiency more effectively and with less bias stands as a significant open problem.

\end{document}